\chapter{关键推导补全}

正文各章按原作者的叙述节奏展开，若干结论给出了结果而略去了中间步骤。本附录把这些跳步
集中补齐，便于读者（尤其是偏好自己复算的读者）验证。每节开头都注明其对应正文的哪一章、
以及补充的是哪一处推导；已在前文完整给出的步骤此处不再重复。

\section{质量频率 \texorpdfstring{$\wms$}{omega\_m} 的引入}
\emph{补充第 1 章与《几何与运动》中“质量频率”定义的来源。}

在洛伦兹宇宙中，自由粒子的能量--动量关系为（取 $c=1$）
\begin{equation}
E^{2}-p^{2}=m^{2},
\end{equation}
即四维动量 $P=(p_x,p_y,p_z,E)$ 在洛伦兹度规下的自内积等于 $m^{2}$。在正交宇宙中，度规
变为全正号，同一物理内容的关系式相应地变为
\begin{equation}
E^{2}+p^{2}=m^{2},
\tag{B.1}
\end{equation}
即 $P\cdot P=m^{2}$（黎曼内积）。空间中静止的粒子对应 $p=0$，此时 $E=m$。

在量子力学中，能量与角频率由 de Broglie 关系 $E=\hbar\omega$ 联系。把静止粒子的能量
$E=m$ 代入，得 $\hbar\omega=m$，即
\begin{equation}
\wms=\frac{m}{\hbar}.
\end{equation}
还原量纲（把 $c$ 放回）即
\begin{equation}
\wms=\frac{mc^{2}}{\hbar},
\end{equation}
这就是\emph{质量频率}（mass frequency）：它是以频率为量纲的粒子静止质量。把 (B.1) 除以
$\hbar^{2}$ 并写作 $\omega^{2}+|\mathbf{k}|^{2}=\wms^{2}$（其中 $p=\hbar|\mathbf{k}|$），
正是正文各章反复使用的色散关系。

\section{黎曼 Proca 方程与电荷守恒}
\emph{补充《黎曼电磁学》中由 (RVWS) 与横波条件导出 Proca 方程、以及导出 (1) 的步骤。}

从带源项的黎曼矢量波方程出发：
\begin{align}
\partial_{x}^{2}\bA+\partial_{y}^{2}\bA+\partial_{z}^{2}\bA+\partial_{t}^{2}\bA+\wms^{2}\bA+\bj&=\mathbf{0}, \tag{B.2}\\
\partial_{x}A^{x}+\partial_{y}A^{y}+\partial_{z}A^{z}+\partial_{t}A^{t}&=0. \tag{B.3}
\end{align}
在黎曼空间的直角坐标中，度规为单位矩阵，故 $\partial_{b}\partial^{b}=\partial_{x}^{2}+\partial_{y}^{2}
+\partial_{z}^{2}+\partial_{t}^{2}$。用 $A^{a}$ 的某个分量记 (B.2) 左端，可写为
\begin{equation}
\partial_{b}\partial^{b}A^{a}=-\wms^{2}A^{a}-j^{a}.
\tag{B.4}
\end{equation}

下一步把电磁场张量 $F_{ab}=\partial_{a}A_{b}-\partial_{b}A_{a}$ 的散度与 (B.4) 联系
起来。直接计算：
\begin{align}
\partial_{b}F^{ab}
&=\partial_{b}\bigl(\partial^{a}A^{b}-\partial^{b}A^{a}\bigr)\nonumber\\
&=\partial^{a}\bigl(\partial_{b}A^{b}\bigr)-\partial_{b}\partial^{b}A^{a}\nonumber\\
&=-\partial_{b}\partial^{b}A^{a},
\tag{B.5}
\end{align}
其中最后一步用到了横波条件 (B.3)：$\partial_{b}A^{b}=0$。把 (B.4) 代入 (B.5)，
\begin{equation}
\partial_{b}F^{ab}=\wms^{2}A^{a}+j^{a},
\end{equation}
即黎曼 Proca 方程
\begin{equation}
\partial_{b}F^{ab}-\wms^{2}A^{a}-j^{a}=0 . \tag{B.6}
\end{equation}
与洛伦兹麦克斯韦方程 $\partial_{b}F^{ab}-j^{a}=0$ 相比，只多出质量项 $-\wms^{2}A^{a}$；
正是它把库仑势由指数衰减改为空间振荡（见正文《黎曼电磁学》）。

再对 (B.6) 取散度 $\partial_{a}$。由于 $F^{ab}$ 反对称，$\partial_{a}\partial_{b}F^{ab}=0$
恒成立（对 $a,b$ 交换后求和相消），于是
\begin{equation}
0=\wms^{2}\,\partial_{a}A^{a}+\partial_{a}j^{a}=\partial_{a}j^{a},
\end{equation}
其中又用了 (B.3)。这就得到正文的方程 (1)：
\begin{equation}
\partial_{x}j^{x}+\partial_{y}j^{y}+\partial_{z}j^{z}+\partial_{t}j^{t}=0,
\end{equation}
即\textbf{电荷守恒}（conservation of charge）：电荷密度随时间的变化率等于电流密度的负
散度。

\section{黎曼 Dirac 方程的四元数形式}
\emph{补充《黎曼量子力学》中将平面波 Dirac 方程约化为四元数方程的一步。}

在 Weyl 基下，黎曼 Dirac 方程的平面波解写作一对左手、右手的二分旋量 $(\xi_{L},\xi_{R})$，
它们与能量--动量矢量 $k=(k_x,k_y,k_z,\omega)$ 通过
\begin{equation}
k^{\mu}H_{\mu}\,\xi_{R}-m\,\xi_{L}=0
\tag{B.7}
\end{equation}
相联系，其中 $H_{\mu}$（$\mu=x,y,z,t$）是张开 $2\times2$ 复矩阵空间 $\mathbb{H}$ 的四个
基矩阵，其乘法表与四元数基 $i,j,k,1$ 一致。

若把 $\mathbb{H}$ 与 $\mathbb{C}^{2}$ 都等同于四元数，则 (B.7) 中出现的三个对象分别按
$\mathrm{SU}(2)\times\mathrm{SU}(2)$ 的不同表示变换。记
\begin{equation}
Q_{V}=k^{\mu}H_{\mu},\qquad q_{L}\leftrightarrow \xi_{L},\qquad q_{R}\leftrightarrow\xi_{R},
\end{equation}
则在元素 $(g,h)\in \mathrm{SU}(2)\times\mathrm{SU}(2)$ 下
\begin{align}
\rho_{(\frac12,\frac12)}(g,h)\,Q_{V}&=g\,Q_{V}\,h^{-1},\\
\rho_{(\frac12,0)}(g,h)\,q_{L}&=g\,q_{L},\\
\rho_{(0,\frac12)}(g,h)\,q_{R}&=h\,q_{R}.
\end{align}
于是乘积 $Q_{V}q_{R}$ 的变换为
\begin{equation}
Q_{V}q_{R}\ \longmapsto\ (g\,Q_{V}\,h^{-1})(h\,q_{R})=g\,(Q_{V}\,q_{R}),
\end{equation}
恰好与 $q_{L}$ 的变换方式 $q_{L}\mapsto g\,q_{L}$ 相同。这就解释了为什么 (B.7) 是一个
\emph{不变}方程：方程两端属于同一个表示 $\rho_{(\frac12,0)}$。据此，把旋量认同为四元数后，
(B.7) 便是极其简洁的三四元数线性关系
\begin{equation}
Q_{V}\,q_{R}-m\,q_{L}=0 . \tag{B.8}
\end{equation}
这正是正文给出的“四元数形式的黎曼 Dirac 方程”。

\section{黎曼 Schwarzschild 解与径向距离}
\emph{补充《黎曼广义相对论》中 Schwarzschild--Kottler 度规的求解思路与 $\operatorname{dist}(r)$ 的积分。}

按球对称、静态的要求设度规
\begin{equation}
\dd s^{2}=A(r)\,\dd t^{2}+B(r)\,\dd r^{2}+r^{2}\bigl(\dd\theta^{2}+\sin^{2}\theta\,\dd\phi^{2}\bigr).
\end{equation}
在黎曼情形下，$t$ 方向与空间方向同号，故 $A,B>0$。把此度规代入真空 Einstein 方程
$G+\Lambda g=0$，得到关于 $A,B$ 的两个方程：其一是
\begin{equation}
\frac{1}{r^{2}}\left(1-\frac{1}{B}\right)-\frac{B'}{rB^{2}}+\Lambda=0,
\end{equation}
其解为
\begin{equation}
\frac{1}{B(r)}=1+\frac{2M}{r}-\frac{\Lambda r^{2}}{3}.
\end{equation}
其二给出 $A(r)=1/B(r)$，即
\begin{equation}
A(r)=1+\frac{2M}{r}-\frac{\Lambda r^{2}}{3}.
\end{equation}
这正是正文给出的黎曼 Schwarzschild--Kottler 度规（洛伦兹版仅 $2M/r$ 项前多一个负号）。

特别地，当 $\Lambda=0$ 时，$t,\theta,\phi$ 固定、只有 $r$ 变化的路径的固有长度为
\begin{equation}
\dd(\text{dist})=\sqrt{1+\frac{2M}{r}}\,\dd r ,
\end{equation}
于是从坐标中心到半径 $r$ 的距离为
\begin{equation}
\operatorname{dist}(r)=\int_{0}^{r}\frac{\dd r'}{\sqrt{1+2M/r'}} .
\tag{B.9}
\end{equation}
做代换 $r'=2M\sinh^{2}u$，则 $\dd r'=4M\sinh u\cosh u\,\dd u$，
$\sqrt{1+2M/r'}=\cosh u/\sinh u$，被积函数化为
\begin{equation}
\frac{\dd r'}{\sqrt{1+2M/r'}}=4M\sinh^{2}u\,\dd u=2M(\cosh 2u-1)\,\dd u .
\end{equation}
积分得
\begin{equation}
\operatorname{dist}(r)=M\sinh 2u-2M u=2M\sinh u\cosh u-2M u .
\end{equation}
由 $\sinh u=\sqrt{r/2M}$、$\cosh u=\sqrt{1+r/2M}$、$u=\operatorname{arcsinh}\sqrt{r/2M}$，
回代即得正文结果
\begin{equation}
\operatorname{dist}(r)=\sqrt{r(r+2M)}-2M\,\operatorname{arcsinh}\!\sqrt{\frac{r}{2M}} .
\end{equation}
由于 $\operatorname{dist}(r)<r$，半径 $r$ 的球面“看起来”比欧氏几何中更大，这正是负曲率
空间的标志，也是黎曼 Schwarzschild 解中不存在视界（故没有黑洞）的几何根源。

\section{黎曼黑体谱：态密度与积分}
\emph{补充《黎曼热力学》中态密度 $\rho(\omega)$、经典总能量密度与黑体的计算。}

在边长为 $L$、体积 $V=L^{3}$ 的立方腔中，周期性边界条件给出平面波波矢
\begin{equation}
\bk_{i,j,k}=\frac{2\pi}{L}(i\,\be_{x}+j\,\be_{y}+k\,\be_{z})+\omega_{i,j,k}\,\be_{t},
\qquad
\omega_{i,j,k}=\sqrt{\wms^{2}-\frac{4\pi^{2}}{L^{2}}(i^{2}+j^{2}+k^{2})} .
\end{equation}
色散关系给出 $\omega^{2}=\wms^{2}-\dfrac{4\pi^{2}}{L^{2}}r^{2}$，其中 $r^{2}=i^{2}+j^{2}+k^{2}$。
频率落在 $[\omega,\omega+\dd\omega]$ 内的模式数，等于（三种偏振乘以）球壳体积
$4\pi r^{2}|\dd r|$。由 $r^{2}=\dfrac{L^{2}}{4\pi^{2}}(\wms^{2}-\omega^{2})$ 对 $\omega$
微分，得
\begin{equation}
2r\,\dd r=-\frac{L^{2}}{2\pi^{2}}\,\omega\,\dd\omega
\quad\Longrightarrow\quad
r\,|\dd r|=\frac{L^{2}}{4\pi^{2}}\,\omega\,\dd\omega .
\end{equation}
于是
\begin{equation}
4\pi r^{2}|\dd r|=\frac{4\pi r^{2}}{r}\bigl(r\,|\dd r|\bigr)
=4\pi r\cdot\frac{L^{2}\omega}{4\pi^{2}}\,\dd\omega
=\frac{L^{2}r\,\omega}{\pi}\,\dd\omega .
\end{equation}
代入 $r=\dfrac{L}{2\pi}\sqrt{\wms^{2}-\omega^{2}}$、$V=L^{3}$，得到
\begin{equation}
4\pi r^{2}|\dd r| = \frac{L^{3}}{2\pi^{2}}\,\omega\sqrt{\wms^{2}-\omega^{2}}\,\dd\omega
= V\,\frac{\omega\sqrt{\wms^{2}-\omega^{2}}}{2\pi^{2}}\,\dd\omega .
\end{equation}
再乘偏振因子 $3$，即得正文的态密度
\begin{equation}
\rho(\omega)\,\dd\omega=3V\,\frac{\omega\sqrt{\wms^{2}-\omega^{2}}}{2\pi^{2}}\,\dd\omega .
\tag{B.10}
\end{equation}

\textbf{经典总能量密度。} 能均分定理给出每个模式平均能量 $kT$，故
\begin{equation}
u_{classical}(\omega)\,\dd\omega=3kT\,\frac{\omega\sqrt{\wms^{2}-\omega^{2}}}{2\pi^{2}}\,\dd\omega .
\end{equation}
对全频段积分。做代换 $\omega=\wms\sin\theta$：
\begin{equation}
\int_{0}^{\wms}\omega\sqrt{\wms^{2}-\omega^{2}}\,\dd\omega
=\wms^{3}\int_{0}^{\pi/2}\sin\theta\cos^{2}\theta\,\dd\theta
=\frac{\wms^{3}}{3},
\end{equation}
因此
\begin{equation}
u_{classical,\,total}=\frac{3kT}{2\pi^{2}}\cdot\frac{\wms^{3}}{3}=\frac{\wms^{3}kT}{2\pi^{2}},
\end{equation}
与正文一致：黎曼经典总能量密度按温度\emph{一次}标度，而非洛伦兹情形的 $T^{4}$。

\textbf{发射功率谱。} 腔内辐射并非各向同性地逃逸：单位时间、单位面积、给定取向的孔所
接收到的辐射，要乘速度 $v(\omega)=\sqrt{\wms^{2}-\omega^{2}}/\omega$，并对方向积分得因子
$1/4$。量子化后每模平均能量为 $\hbar\omega/(\ee^{\hbar\omega/kT}-1)$，与 (B.10) 相乘并
乘 $v(\omega)/4$，即得正文的发射谱
\begin{equation}
I_{quantum}(\omega)\,\dd\omega=\frac{3\hbar\omega\,(\wms^{2}-\omega^{2})}{8\pi^{2}\bigl(\ee^{\hbar\omega/kT}-1\bigr)}\,\dd\omega .
\end{equation}

\section{黎曼矢量波方程的三偏振}
\emph{补充《几何与波》中“横波条件等价于剔除标量波”的论证。}

考虑矢量平面波 $\bA(x)=\bA_{0}\sin(\bk\cdot\bx)$，其中 $\bk\cdot\bx=k_x x+k_y y+k_z z+\omega t$。
只要 $\bk\cdot\bk=\wms^{2}$，$\bA$ 的每个分量都满足黎曼标量波方程（RSW）。然而，可以由此
构造一个所有观测者都认同的\emph{标量}波
\begin{equation}
D(x)=\bk\cdot\bA(x)=(\bk\cdot\bA_{0})\sin(\bk\cdot\bx).
\end{equation}
为使 $\bA$ 表示纯矢量波（对应自旋 $1$ 的量子，如黎曼光子）而不含标量成分（自旋 $0$），
须令
\begin{equation}
\bk\cdot\bA_{0}=0 .
\tag{B.11}
\end{equation}
这一条件把 $\bA_{0}$ 限制在垂直于传播矢量 $\bk$ 的三维子空间内，故黎曼光有\emph{三}个
独立偏振（对比我们宇宙的光只有两个）。

把 (B.11) 写成对任意场 $\bA(x)$ 都适用的局域形式。对上述平面波，
\begin{align}
\partial_{x}A^{x}+\partial_{y}A^{y}+\partial_{z}A^{z}+\partial_{t}A^{t}
&=(A_{0}^{x}k^{x}+A_{0}^{y}k^{y}+A_{0}^{z}k^{z}+A_{0}^{t}k^{t})\cos(\bk\cdot\bx)\nonumber\\
&=(\bk\cdot\bA_{0})\cos(\bk\cdot\bx)=0,
\end{align}
即\emph{横波条件}
\begin{equation}
\partial_{x}A^{x}+\partial_{y}A^{y}+\partial_{z}A^{z}+\partial_{t}A^{t}=0 .
\tag{B.12}
\end{equation}
(B.12) 与 RSW 一起构成正文的黎曼矢量波方程（RVW）。三个偏振在量子化后即黎曼光子的三个
自旋态，这一点在《黎曼量子力学》与《黎曼热力学》的态计数中反复使用。
